\chapter{系统实现}
\label{chap:impl}

\section{模块组织}

实现基于 xv6-riscv。主要模块包括：VFS 路径遍历与 \texttt{/pm} 挂载点；\texttt{pmfile}（普通 fd、COW 元数据、硬链接/rename）；\texttt{pmcdc}/\texttt{pmstream}；\texttt{pmalgo}/\texttt{pmindex}；\texttt{pmfs}（checkpoint、可达性、压缩）；\texttt{pmdev}/\texttt{pmem}；以及 \texttt{tools/pmfs\_recovery.py}、\texttt{test-xv6.py}、契约与 bench 脚本。文件内容经 \texttt{pm\_put}/\texttt{seal} 落入 run；元数据经 COW 页与 checkpoint 发布。

\section{VFS 与普通文件路径}

\texttt{vfs\_walk} 逐组件解析；仅原生根下的 \texttt{pm} 进入 PMFS，\texttt{/pm/..} 返回原生根。PM inode pin 支持 cwd/目录 fd；独立 \texttt{open} 各自维护 offset，共享工作缓冲；\texttt{dup}/\texttt{fork} 共享打开对象。PMFS sleep lock 串行化缓冲、append 取尾与同步。

硬链接维护共享 inode 与链接数；非最后链接可在打开时删除，最后链接的 unlink-open 仍拒绝。子树 rename 检查环与路径长度后，修改内存图并以\textbf{一次} checkpoint 发布；非替换可复用不可变 inode 页。

\section{COW 元数据、run 与 seal}

内容同步：CDC/固定切分 $\rightarrow$ \texttt{pm\_put} $\rightarrow$ seal 追加不可变 run $\rightarrow$ 新 inode 快照（可复用未变 dentry）$\rightarrow$ checkpoint。\texttt{pm\_get} 优先 memtable/缓存，再精确索引或 newest-first + Bloom/fence。索引满后未跟踪键走 fallback；更新已跟踪键必须刷新指向。

安全挂载：无有效/旧格式/损坏则失败关闭；仅 \texttt{pmformat --force} 显式格式化。v2 与 v3 不兼容，无迁移工具。

\section{刷写、设备与可选 CBO}

\texttt{pmdev\_write\_page} 后调用 \texttt{pmem\_persist}。默认 flush 逐 64\,B 行 touch 并计数；\texttt{pmem\_fence} 含编译器屏障与内存 fence。可选 \texttt{PMEM\_ZICBOM} 发出 \texttt{CBO.CLEAN}（假定块大小与权限）。实验性 NVDIMM：补丁 QEMU 注册独立设备，内核在 \texttt{kalloc} 前解析私有 FDT 并校验地址/大小。

\section{维护：页复用与显式压缩}

可达性扫描覆盖两槽根与 run 链；不可达页可分配给新元数据或新 run。\texttt{pmcompact} 要求 memtable/stream 空，预留不相交暂存后写出新链并一次发布；忙或空间不足返回明确错误。这不是在线全设备 GC。

\section{验收脚手架}

\begin{description}[leftmargin=1.5em,style=nextline]
  \item[旧自测] \texttt{ONBOX\_PMFS\_*} 阶段码；\texttt{contract-full} 等破坏性，仅 disposable。
  \item[重启 oracle] \texttt{make test-recovery*}：临时后端、guest 步骤、\texttt{SIGKILL}、同后端读回；覆盖挂载拒绝、文件/KV、目录/链接、插桩切点。
  \item[模型] \texttt{make test-model}：有界协议枚举 + 负控反例。
  \item[轨迹] 测试构建记录写范围/flush/fence，采样重放提前逐出等。
\end{description}

尖刺价值：正确性依赖「可脚本判定 + 诚实假设」，而非口头演示。
